Der Stein macht $n = \nStO$ Umdrehungen in $t = \tStO$, die Frequenz ist also
\begin{align}
 f &= \fStF = \fSt
\end{align}
Die Schnur liefert die Zentripetalkraft. Mit $\omega = 2\pi f$ gilt
\begin{align}
 F &= m\,\omega^2 r = 4\pi^2 f^2\, r\, m
\end{align}
Aufgelöst nach der Masse (die Gewichtskraft des Steins und die Neigung der Schnur werden vernachlässigt, die Schnurlänge ist der Kreisradius):
\begin{align}
 m &= \mStF \\
   &= \frac{\FSt}{4\pi^2\cdot\left(\fSt\right)^2\cdot\rSt} \\
   &= \mSt \approx \result{\mStgP{0}{2}}
\end{align}
